\section{数据与测度}
\sectionlead{section-2}{二、数据与测度}{怎样从网页链接记录识别投资者研究的标的、内容和订单时点？}

\subsection{2.1 浏览样本与 URL}
\begin{frame}[label=sample]{网页浏览数据描述}
  \crumb{Part 2}{2.1 浏览样本与 URL}{II.A · Table 1}
  \pair{\textbf{原始在线家庭面板}\par\medskip
    \greatnum{20 亿}\;条 URL\par\smallskip
    \textbf{数万户家庭}\par\small
    美国，2007 年 1、2、3、6 月。\par
    观察月份因户而异；随机拨号等方式招募。}
  {\textbf{至少有一笔推定股票订单}\par\medskip
    \greatnum{484}\;户\par\smallskip
    \textbf{850 万点击}\quad\textbf{6 万小时}\par\smallskip
    \textbf{2,911 笔}推定普通股／ADR 订单}
  \smallskip
  {\color{wine}\rule{.94\linewidth}{.5pt}}\par\smallskip
  {\small 匿名家庭 ID \(\mid\) 时间戳 \(\mid\) URL \(\mid\) 停留时长\par\smallskip
  8 家券商，7 家识别到订单；站外含 Yahoo Finance、CNBC、MarketWatch。}
  \finding{\textbf{数据来自}数万户家庭的浏览器插件记录。\quad
  \textbf{缺少}后台成交确认与完整持仓。}
\end{frame}

\newcommand{\udomain}[1]{\textcolor{quiet}{#1}}
\newcommand{\upath}[1]{\textcolor{sampleteal}{#1}}
\newcommand{\ucode}[1]{\textcolor{sampleblue}{#1}}
\newcommand{\uparam}[1]{\textcolor{samplegold}{#1}}
\newcommand{\utrade}[1]{\textcolor{wine}{#1}}
\begin{frame}[label=urls]{三类 URL 如何读出股票、主题与展示参数}
  \crumb{Part 2}{2.1 浏览样本与 URL}{Table IA1 第 20、13、7 行 · 原文构造例}
  {\small\textbf{1 下单}\quad\utrade{卖出 100 股}\quad\ucode{GOOG}\par}
  {\fontsize{8}{10}\selectfont\ttfamily
  \udomain{https://[broker-1].com}\upath{/invest/socreateentry?}\par
  \utrade{ordertype="sell"}\ucode{\&symbol=GOOG}\utrade{\&shares=100}\par}
  \smallskip
  {\small\textbf{2 研究}\quad\upath{分析师意见}\quad\ucode{IMCL}\par}
  {\fontsize{8}{10}\selectfont\ttfamily
  \udomain{http://finance.yahoo.com}\upath{/q/ao?}\ucode{s=IMCL}\par}
  \medskip
  {\small\textbf{3 图表}\quad\upath{盈利图}\quad\uparam{回看 1,095 天、前瞻 545 天}\par}
  {\fontsize{8}{10}\selectfont\ttfamily
  \udomain{https://[broker-1].com}\upath{/stocks/earnings/update\_earningsChart.asp?}\par
  \uparam{showQTRChange=1\&display=mountain\&overlay=}\par
  \uparam{\&indicator=volume,1,13,EMA,13\&duration=1095\&daysforward=545\&year=2006}\par}
  \smallskip{\small IMCL 由相邻链接定位；参数还显示图形格式及成交量 EMA。}\par\medskip
  {\footnotesize\udomain{域名：来源}\quad\upath{路径：主题}\quad\ucode{代码：股票}\quad
   \uparam{参数：展示}\quad\utrade{下单字段：方向与数量}}
\end{frame}

\begin{frame}[label=yahoo]{Yahoo Finance 的分析师页面保留相似结构}
  \crumb{Part 2}{2.1 浏览样本与 URL}{Figure A1}
  {\centering
  \annotatedpanel{4.10cm}{a1.pdf}{
    \draw[wine,semithick,-{Stealth[length=1.2mm]}] (.438,.943)--(.561,.855);
    \draw[wine,semithick,-{Stealth[length=1.2mm]}] (.438,.773)--(.561,.982);
    \draw[sampleteal,semithick,-{Stealth[length=1.2mm]}] (.438,.593)--(.561,.598);
    \draw[samplegold,semithick,-{Stealth[length=1.2mm]}] (.438,.358)--(.561,.258);
  }\par}
  {\small\textbf{2007 年 2 月 23 日}\hfill\textbf{2026 年 4 月 15 日}\par}\smallskip
  \pair{\small\textbf{网站页面}\quad 模块变化不大，Yahoo 仍是重要信息平台。}
       {\small\textbf{用户行为}\quad 移动端、社交媒体与 AI 咨询兴起，获取信息的方式已改变。}
\end{frame}

\subsection{2.2 记录整理与交易归属}
\begin{frame}[label=rules]{浏览会话与重复订单的整理}
  \crumb{Part 2}{2.2 记录整理与交易归属}{Table IA1 · Steps 1、2 · 原文构造例}
  {\small 2007 年 3 月 23 日，22 条 URL 串起约 14 分钟的浏览。}\par\bigskip
  \begin{center}
  \begin{tikzpicture}[x=.70cm,y=1cm,font=\scriptsize]
    \foreach \a/\b/\c in {0/5.15/papergray,5.15/6.32/paperblue,6.32/8.05/paperblue,8.05/9.68/paperblue,9.68/11.87/papergray,11.87/14.33/paperrose}
      \fill[\c] (\a,0) rectangle (\b,.60);
    \foreach \x/\n in {0/2:31:11,5.15/2:36:20,8.05/2:39:14,11.87/2:43:03,14.33/2:45:31}
      \draw[quiet] (\x,0)--(\x,-.08) node[below]{\n};
    \node at (2.58,.30) {赛车};\node at (5.73,.30) {市场};
    \node[align=center] at (7.18,.30) {券商\\IMCL};
    \node[align=center] at (8.87,.30) {Yahoo\\IMCL};
    \node at (10.78,.30) {赛车};\node at (12.7,.30) {GOOG};
    \draw[wine,thick] (13.85,-.03)--(13.85,.88) node[above left]{2:45:02 卖出 100 股};
  \end{tikzpicture}
  \end{center}\medskip
  \leadrow{会话划分}{相邻点击间隔 \(\leq15\) 分钟，归入同一浏览会话。}
  \medskip
  \leadrow{订单合并}{同股、同向候选订单相隔 \(\leq30\) 分钟则合并，取末次时点。}
  \finding{这次会话研究过 IMCL，最后下单的是 GOOG。全部研究与交易标的研究需要区分。}
\end{frame}

\begin{frame}[label=assignment]{补齐代码，再把研究归到同股订单}
  \crumb{Part 2}{2.2 记录整理与交易归属}{Step 8 · p.\,37 · 归属教学例}
  \small\textbf{补缺码：同会话、前后最近已匹配记录为同一股}\par\smallskip
  {\centering
  \begin{tikzpicture}[node distance=7mm,font=\small]
    \node[storybox,minimum height=.7cm,text width=2.2cm] (a) {IMCL 研究};
    \node[storyfocus,minimum height=.7cm,text width=2.2cm,right=of a] (b) {缺码新闻};
    \node[storybox,minimum height=.7cm,text width=2.2cm,right=of b] (c) {IMCL 研究};
    \draw[storyarrow] (a)--(b);\draw[storyarrow] (b)--(c);
    \draw[wine,dashed,-{Stealth}] (a.north) to[bend left=25] node[above,font=\scriptsize]{补配 IMCL} (b.north);
  \end{tikzpicture}\par}\smallskip
  \small\textbf{归研究：对准下一笔同股交易（同一股票的教学例）}\par\smallskip
  {\centering
  \begin{tikzpicture}[node distance=4mm,font=\small]
    \node[storybox,text width=2.0cm,minimum height=.85cm] (t) {周二\\研究};
    \node[storyfocus,text width=2.0cm,minimum height=.85cm,right=of t] (w) {周三\\卖出};
    \node[storybox,text width=2.0cm,minimum height=.85cm,right=of w] (h) {周四\\研究};
    \node[storyfocus,text width=2.0cm,minimum height=.85cm,right=of h] (f) {周五\\买入};
    \draw[grapharrow] (t)--(w);\draw[grapharrow] (h)--(f);
  \end{tikzpicture}\par}\smallskip
  {\small 不是固定归到次日；之后无同股交易，则归到末笔交易之后。}
  \finding{得到每位投资者、每只股票、带主题和时间的研究记录，以及可归属的订单时点。}
  \scriptsize\jump{appendix-4}{A4 · 原文处理步骤}
\end{frame}
